
%=======================================================================================
%=======================================================================================
\chapter{Data}
This chapter describes the dataset used for model training and evaluation, outlining key feature engineering steps and data limitations. Table~\ref{tab:data_summary} summarises all predictor and target variables alongside their sources and temporal coverage. Feature selection rationale, data partitioning, and specific modelling choices are detailed in Chapter~\ref{ch:methodology}.

\section{Study area and cohort coverage}
The dataset covers approximately $16\text{~km} \times 12\text{~km}$ of Aberfoyle and contains 795,645 plot--survey records from 230,359 unique plot locations across six survey years (2002--2023). The records contain airborne laser scanning (ALS)-derived variables, including canopy cover and height percentiles, alongside stand and management fields such as planting year and known felling history (Table~\ref{tab:data-fields}). Figure~\ref{fig:data-study-area} shows the study-area topography, the spatial coverage of both cohorts, and an example of standard and boundary-clipped plots.

\begin{figure}[!htbp]
\includegraphics[
    width=1.05\textwidth,
    height=0.48\textheight,
    keepaspectratio
]{figures/fig_data/aberfoyle_area.png}
\caption{\small{Left:Aberfoyle study area. Middle: cohort coverage, Right: plot-level detail}}
\label{fig:data-study-area}
\end{figure}

Because LiDAR coverage increase by $37\text{ km}^2$ between 2006 and 2008, two nested cohorts were defined (Figure~\ref{fig:cohort_coverage_funnel}): a broad-coverage \emph{4-survey} cohort (2008, 2012, 2021, 2023; 71,766 plots across 296 compartments) and a longer-trajectory \emph{6-survey} cohort (2002--2023; 13,897 plots across 48 compartments).

To ensure data quality, sequential filters were applied to both cohorts (Table~\ref{tab:data-cleaning-funnel}):
\begin{itemize}
    \item \textbf{Complete presence:} Plots must be present across all survey years in the cohort.
    \item \textbf{Species purity:} Pure Sitka spruce stands only.
    \item \textbf{Plausible records:} Planting year between 1850 and the survey year, stand age $0\text{--}200$ years, and LiDAR top height $0\text{--}60\text{ m}$ (ecological maximum \cite{mason2011Sitka_C1_SS}).
    \item \textbf{Maturity cutoff:} Stand age $\ge 30$ years by 2023 to exclude unreliable young-stand height measurements \cite{suarezminguez2008Practical_C1_SS}.
\end{itemize}

\begin{table}[!htbp]
\centering
\scriptsize
\begin{tabular}{lrrrr}
\toprule
& \multicolumn{2}{c}{4-survey} & \multicolumn{2}{c}{6-survey} \\
\cmidrule(lr){2-3}\cmidrule(lr){4-5}
Stage & Plots & Removed & Plots & Removed \\
\midrule
0. Raw dataset
& 230,359 & -- & 230,359 & -- \\
1. Cohort years present
& 230,345 & 0.01\% & 230,359 & 0.00\% \\
2. Complete survey coverage
& 112,782 & 51.04\% & 18,781 & 91.85\% \\
3. Sitka spruce only
& 72,395 & 35.81\% & 13,897 & 26.01\% \\
4. Valid planting year
& 71,766 & 0.87\% & 13,897 & 0.00\% \\
5. Plausible age (0--200 years)
& 71,766 & 0.00\% & 13,897 & 0.00\% \\
6. Plausible height (0--60 m)
& 71,766 & 0.00\% & 13,897 & 0.00\% \\
7. Mature stands (age $\geq 30$ in 2023)
& \textbf{58,112} & 19.03\%
& \textbf{13,769} & 0.92\% \\
\bottomrule
\end{tabular}
\caption{\small{Plots left after each filter. Percentages show the share removed from the previous row, rather than from the raw dataset. Figure~\ref{fig:data-study-area} shows the final spatial coverage.}}
\label{tab:data-cleaning-funnel}
\end{table}

%------------------------------------------------
\section{Response variable and derived attribution targets}

All targets in this dissertation are built from the plot-level 95th percentile ALS return height (\texttt{elev\_percentile\_95th}). From this base metric, three distinct targets are derived (detailed in Chapter~\ref{ch:methodology}): the mean residual departure from a shared regional growth curve across all survey years (Q1), the difference between a plot's fitted asymptote and its assigned yield-class asymptote (Q2), and the raw LiDAR top height itself (Q3). Forest Research records also supply planting year, management history, and yield estimates (Table~\ref{tab:data-fields}), with any variables calculated directly from measured height removed from model inputs to prevent target data leakage. A code-level audit confirmed that the Q2 target contains no circularity in this project's pipeline. The target was calculated strictly from Forest Research inventory records, while environmental variables (terrain, climate, wind, and soil) were extracted independently from external GIS layers and joined only at the modeling stage. This prevents internal data leakage, though it cannot verify whether historical expert judgement originally influenced Forest Research yield-class assignments.


\begin{table}[!htbp]
\scriptsize
\setlength{\tabcolsep}{3pt}
\renewcommand{\arraystretch}{0.95}
\begin{tabular}{
    >{\raggedright\arraybackslash}p{2.6cm}
    >{\raggedright\arraybackslash}p{7.0cm}
    >{\raggedright\arraybackslash}p{3.2cm}
}
\hline
Field & Meaning or derivation & Role in this study \\
\hline
\texttt{elev\_\allowbreak percentile\_\allowbreak 95th} & 95th percentile of ALS return heights above ground, used without correction & Response variable \\
\texttt{Top\_Height95} & $=\texttt{elev\_percentile\_95th}\times1.1$; Forest Research's own corrected height estimate & Retained for reporting but not modelled \\
\texttt{plyr}, \texttt{Age} & Planting year; \texttt{Age} = survey year $-$ \texttt{plyr} & Input feature; used for data cleaning \\
\texttt{yldc} & An existing Forest Research yield-class field, not calculated from this project's LiDAR-derived height response & Excluded from prediction inputs; used to construct Q2's own attribution target \\
\texttt{whcl} & Ordinal Forest Research windthrow hazard class, 0--6 & Management descriptor \\
\texttt{CanopyCover} & Canopy-cover fraction from the same ALS survey and processing step as the height response & Input feature \\
\texttt{Thin}, \texttt{last\_thinn}, \texttt{next\_thin\_} & Thinning history; \texttt{Thin} is binary (1 if thinned, 0 if never thinned); \texttt{last\_thinn}/\texttt{next\_thin\_} record last and next thinning years & Input feature \\
\texttt{Vol95}, \texttt{area}, \texttt{g1}, \texttt{g3} & Forest Research volume estimate, calculated from height, plot area, and canopy cover & Used for disturbance screening \\
\texttt{GYCspec95} & General Yield Class calculated directly from a plot's own measured height (\texttt{Top\_Height95}) & Excluded (target leakage) \\
\texttt{p1}--\texttt{p5} & Forest Research curve parameters associated with yield class & Excluded from inputs; \texttt{p4}/\texttt{p5} held fixed to build Q2 target \\
\hline
\end{tabular}
\caption{\small{Target, stand, and management fields provided by Forest Research. \texttt{GYCspec95} and \texttt{p1}--\texttt{p5} are excluded from model inputs to prevent target data leakage. In Q2, \texttt{p4} and \texttt{p5} are held fixed only to define the baseline yield-class asymptote target.}}
\label{tab:data-fields}
\end{table}


%------------------------------------------------
\section{Derived terrain, wind, climate and soil variables}
\label{sec:data-environmental}

Each plot was matched with external terrain, wind, climate, and soil variables (Table~\ref{tab:data-env-sources}). Appendix~\ref{tab:data-env-variables_detail_app} defines each variable.

\begin{table}[!htbp]
\centering
\scriptsize
\setlength{\tabcolsep}{4pt}
\renewcommand{\arraystretch}{1.05}
\begin{tabular}{ >{\raggedright\arraybackslash}p{2.0cm} >{\raggedright\arraybackslash}p{4.2cm} >{\raggedright\arraybackslash}p{2.3cm} >{\raggedright\arraybackslash}p{4.0cm}}
\hline
Data family & Source and resolution & Time basis & Variables used \\
\hline
Terrain & OS Terrain 50, 50\,m \cite{ordnancesurveyTerrain50_C3_LID} & Static & Elevation, slope, aspect, northness and eastness \\

Soil & CEH Natural Capital, 50\,m \cite{feeney2026CEHNaturalCapital_C3_LID}; SoilGrids v2.0, 250\,m \cite{poggio2021SoilGrids_C3_LID} & Static & Pedotope, drainage, texture and soil pH \\

Wind exposure & Global Wind Atlas, 250\,m \cite{davis2023GlobalWindAtlas_C3_LID} & 2006--2015 climate & Long-term wind speed and exposure \\

Observed wind & MIDAS Open weather stations \cite{metoffice2019MIDASOpen_C3_LID} & Survey intervals & Mean and extreme wind conditions \\

Long-term climate & CHELSA v2.1, 1\,km \cite{karger2017CHELSA_C3_LID} & 1981--2010 climate & Growing degree days and precipitation \\

Matched climate & HadUK-Grid, ~1\,km \cite{hollis2019HadUKGrid_C3_LID} & Cohort years & Temperature and frost \\

Spatial context & OS Open Roads and Open Rivers, vector data
\cite{ordnancesurveyOpenRoads_C3_LID,ordnancesurveyOpenRivers_C3_LID} & Static & Distance to roads, rivers and forest edges \\
\hline
\end{tabular}

\caption{\small{External data sources used in the analysis. Values were matched to plot centres. Full source, resolution, coverage and extraction details are given in Appendix~\ref{app:data-sources}.}}
\label{tab:data-env-sources}
\end{table}

Environmental variables vary in their time resolution. Terrain, soil, climate normals (CHELSA, Global Wind Atlas), and TOPEX were kept static, assigning each plot a single value across all survey years. TOPEX measures topographic shelter, where higher scores mean greater shelter and lower scores mean higher exposure. By contrast, MIDAS wind metrics were calculated specifically for each interval between survey dates using weather-station data. HadUK-Grid temperature and frost values were averaged over the survey years within each cohort, while road, river, and edge distances were treated as static baselines. Limited missing values in environmental features were handled by removing incomplete rows for specific model sets or using missingness indicators for thinning history (Chapter~\ref{ch:methodology}).

Several limitations should be noted. Requiring plots to appear across every survey year ensured clean tracking over time, but reduced landscape coverage, particularly in the 6-survey cohort where high-elevation areas are underrepresented. Predictor layers also vary in detail: terrain changes across individual plots, whereas regional climate maps share the same values across wide areas and do not update per survey date. Additionally, canopy cover shares the same LiDAR flight acquisitions as top height, which could introduce shared survey conditions. These limits are especially relevant when evaluating local plot-level attribution.

